\bookinterlude{I}{II}{Do outro lado da portinhola}
Você se lembra da peça de 5,2 mm que o inspetor mediu duas vezes? Vou atrás dela. A portinhola dá um estalo, a peça cai na gaveta A e eu dou a volta na bancada para ver quem vem esvaziar a gaveta.

É a mulher que faz as buchas em que essas peças vão entrar. Amanhã ela precisa abrir o furo da próxima bucha enquanto a nossa peça ainda está no torno. Pergunta que diâmetro pode esperar. Digo que estou disposto a apostar dinheiro que a peça vai para A. Ela põe o dedo no desenho. “E quantos milímetros eu escrevo aqui?”

O inspetor procura no caderno a regulagem que vai usar amanhã. Nas últimas três vezes saíram duas peças de 5,2 mm e uma de 5,8 mm. Todas foram para A. Você sugere 5,2. Eu faço a média e sugiro 5,4. Naquelas três peças você teria acertado duas em cheio e errado a outra por seis décimos. Eu teria errado as três, duas por dois décimos e uma por quatro.

Ela continua com o lápis na mão. Antes de lhe dar um número, eu queria ouvir a opinião de quem vai fazer os furos. O que dá mais trabalho nessa bancada, corrigir os meus três erros ou o seu?
